\pdfminorversion=7
\documentclass[11pt]{article}

\usepackage[T1]{fontenc}
\usepackage[utf8]{inputenc}
\usepackage{amsmath,amssymb,bm}
\IfFileExists{newtxtext.sty}{\usepackage{newtxtext,newtxmath}}{}
\usepackage{booktabs}
\usepackage{threeparttable}
\usepackage{tabularx}
\usepackage{array}
\usepackage{graphicx}
\usepackage{float}
\usepackage{xcolor}
\usepackage[margin=1in]{geometry}
\definecolor{paperlinkblue}{RGB}{0,51,102}
\definecolor{paperlinkred}{RGB}{150,24,24}
\usepackage[colorlinks=true,linkcolor=paperlinkred,citecolor=paperlinkblue,urlcolor=paperlinkblue]{hyperref}
\urlstyle{rm}
\def\UrlFont{\rmfamily}

\graphicspath{
{../manuscript/figures/round2_candidate/appendix/}
{../manuscript/figures/round2_candidate/main/}
{../manuscript/figures/round2_candidate/manual/}
{../manuscript/figures/}
}
\sloppy
\renewcommand{\thefigure}{S\arabic{figure}}
\renewcommand{\thetable}{S\arabic{table}}
\renewcommand{\figurename}{Fig.}
\expandafter\def\csname nosticsword\endcsname{nostics}
\newcommand{\diagword}{diag\csname nosticsword\endcsname}
\newcommand{\sifigure}[1]{\includegraphics[width=\linewidth,height=0.74\textheight,keepaspectratio]{#1}}
\newcommand{\sitable}[0]{%
  \small
  \setlength{\tabcolsep}{5pt}%
  \renewcommand{\arraystretch}{1.08}%
}

\begin{document}

\title{Supplementary Information for\\
Surrogate-conditioned optimizer ranking in block-scale urban energy design}
\author{Haolin Wang, Zhi Wu, Wei Gu, Xiao Zhang, Pengxiang Liu, Qirun Sun, and Wei Wang}
\date{}
\maketitle

This Supplementary Information contains the case settings, analytic response-generation equations, surrogate assessment results, optimizer comparison details, feasible-projection summaries, physics-based cross-model stress-test outputs, and cross-climate sensitivity tables supporting the main manuscript.

\section*{S1. Case settings and analytic response generation}

\subsection*{S1.1 Block geometry and prototype settings}

Each generated case uses a 240~m $\times$ 240~m block divided into a $3\times3$ grid of 80~m $\times$ 80~m land units. One land unit is assigned as open space, and the remaining eight land units are populated with residential prototypes. The floor-to-floor height is 3.0~m. The prototype set contains point, slab, and courtyard forms with footprint areas of 1000--2000~m$^2$ and prototype-specific floor ranges.

\begin{table}[H]
\caption{Residential prototype settings used by the morphology generator.}
\label{tab:s-prototypes}
\centering
\small
\setlength{\tabcolsep}{6pt}
\renewcommand{\arraystretch}{1.08}
\begin{threeparttable}
\begin{tabular*}{\linewidth}{@{\extracolsep{\fill}}llrrrrr@{}}
\toprule
Prototype & Category & Footprint (m$^2$) & Min floors & Max floors & Width (m) & Depth (m) \\
\midrule
P-1 & Point & 2000 & 1 & 3 & 40 & 50 \\
P-2 & Point & 1500 & 4 & 12 & 30 & 50 \\
P-3 & Point & 1000 & 13 & 30 & 25 & 40 \\
P-4 & Point & 1000 & 13 & 30 & 20 & 50 \\
S-1 & Slab & 2000 & 1 & 3 & 25 & 80 \\
S-2 & Slab & 1500 & 4 & 12 & 20 & 75 \\
S-3 & Slab & 1000 & 13 & 30 & 16 & 62.5 \\
C-1 & Courtyard & 2000 & 1 & 3 & 50 & 50 \\
C-2 & Courtyard & 1500 & 4 & 12 & 45 & 40 \\
\bottomrule
\end{tabular*}
\begin{tablenotes}[flushleft]
\footnotesize
\item The footprint and linear dimensions define the prototype geometry assigned to populated land units before morphology descriptors are calculated.
\end{tablenotes}
\end{threeparttable}
\end{table}

\subsection*{S1.2 Random morphology generation}

For each generated block, the open-space index is sampled from 0 to 8, theta is sampled uniformly from -45~deg to 45~deg, and each populated cell receives a randomly sampled prototype and a floor count sampled within that prototype range. Descriptor values are calculated only after this geometry-generation step. This order keeps the generated feasible block as the primary object and treats the 12 surrogate inputs as derived morphology descriptors.

\subsection*{S1.3 Descriptor definitions and dependencies}

The descriptor set is FAR, SD, AF, AR\_ew, AR\_ns, SVF, BD, OSR, SC, PAR, theta, and OSLI. FAR is numerically linked to BD and AF, while OSR is linked to footprint area and gross floor area. OSLI is integer-valued by construction. In the locked descriptor-dependency check, six principal components explain 95\% of descriptor variance, and exact duplicate descriptor rows are absent in the 2000-row canonical pool.

\begin{table}[H]
\caption{Descriptor dependency checks in the canonical 2000-row pool.}
\label{tab:s-dependencies}
\centering
\footnotesize
\begin{tabularx}{\linewidth}{@{}p{0.35\linewidth}p{0.25\linewidth}p{0.28\linewidth}@{}}
\toprule
Check & Value & Interpretation \\
\midrule
FAR minus BD times AF & $1.33\times10^{-15}$ & algebraic residual is numerical roundoff \\
OSR relation residual & $2.22\times10^{-16}$ & footprint and gross floor area relation holds \\
OSLI integer consistency & 1.0 & all OSLI values are integer \\
Exact duplicate count & 0 & no exact duplicate descriptor rows \\
Nearest-neighbour distance median & 0.199 & generated pool is sparse in normalized descriptor space \\
PCA components for 95\% variance & 6 & descriptor space is structured \\
\bottomrule
\end{tabularx}
\end{table}

\subsection*{S1.4 Analytic EUIt, EG, and H response equations}

The analytic response generator is the target-assignment module that maps a generated morphology descriptor vector, denoted $\mathbf{x}$, to three target responses, denoted $\mathbf{y}^{ana}=[\mathrm{EUIt},\mathrm{EG},H]^\mathsf{T}$. The mapping uses explicit formulas that combine density, openness, aspect-ratio, sky-view, orientation, and open-space-position terms to represent demand, rooftop-generation, and sunlight-access responses. A small stochastic term is added during dataset construction, and the resulting target values are clipped to the training-domain ranges used by the surrogate. This module provides a controlled response surface for optimizer comparison and defines the analytic target regime used in the main manuscript.

Let $\theta$ be the orientation angle, let $OSLI^\ast = OSLI/8$, and define

\begin{align}
p_\theta &= |\theta| / 45,\\
a_m &= 0.5(AR_{ew}+AR_{ns}),\\
s &= \exp\left(-\frac{(FAR-1.7)^2}{0.6}\right)(0.45+SVF)(0.3+OSR),\\
p_d &= \max(FAR-2.8,0).
\end{align}

For the baseline weather condition, the fallback station-bias terms are zero. The unclipped analytic responses are

\begin{align}
EUIt^\ast ={}& 78.0 + 10.0BD + 4.3PAR + 2.8a_m + 0.022SD + 0.8SC + 2.1p_\theta \nonumber\\
& - 6.4OSR - 5.8SVF - 0.9AF - 4.5s + \epsilon_{EUIt},\\
EG^\ast ={}& 0.96 + 0.28FAR + 0.64SVF + 0.52OSR - 0.19BD - 0.12a_m - 0.18p_\theta \nonumber\\
& + 0.26OSLI^\ast + 0.78s - 0.1p_d + \epsilon_{EG},\\
H^\ast ={}& 6.45 + 1.1SVF + 0.55OSR + 0.48OSLI^\ast + 0.24s - 0.7p_\theta - 0.45PAR - 0.24a_m + \epsilon_H.
\end{align}

The Gaussian noise scales are 0.22 for EUIt, 0.03 for EG, and 0.02 for H in the response-generation configuration. The final target values are clipped as

\begin{align}
EUIt &= \operatorname{clip}(EUIt^\ast,66.0,96.0),\\
EG &= \operatorname{clip}(EG^\ast,1.2,2.85),\\
H &= \operatorname{clip}(H^\ast,6.0,7.85).
\end{align}

\subsection*{S1.5 Descriptor distributions and sampling coverage}

Figure~\ref{fig:s1-descriptor-distributions} summarizes the descriptor coverage of the generated morphology pool. The first panel shows the spread of the 12 morphology descriptors, while the OSLI histogram confirms that the open-space position categories are represented across the generated cases. The nearest-neighbour distance distribution provides a compact view of sparsity in normalized descriptor space, which is relevant because the optimizers later search a continuous descriptor domain.

\begin{figure}[H]
\centering
\sifigure{A1_descriptor_distributions.pdf}
\caption{Descriptor distributions and nearest-neighbour distances in the canonical 2000-row generated morphology pool.}
\label{fig:s1-descriptor-distributions}
\end{figure}

\section*{S2. Surrogate selection and robustness assessment}

\subsection*{S2.1 Dataset-scale study}

The nested scale study compares 500, 1000, 1500, and 2000 sample regimes. Each regime evaluates candidate DNN presets, and the selected 2000-sample tuned-standard model gives mean target nMAE = 0.0182, mean tail nMAE = 0.0203, mean $R^2$ = 0.9823, worst target nMAE = 0.0242, and selection objective = 0.0233. This setting is used as the shared surrogate environment for the optimizer comparison.

\begin{figure}[H]
\centering
\sifigure{A3_scale_study.pdf}
\caption{Dataset-scale surrogate-selection comparison across the 500, 1000, 1500, and 2000-row regimes.}
\label{fig:s2-scale-study}
\end{figure}

The scale study shows gradual improvement as the sample pool increases from 500 to 2000 cases. The 2000-sample regime provides the lowest selection objective among the tested regimes and is therefore used for the shared surrogate environment in the main comparison. This selection should be interpreted within the analytic response-generation regime, since the target values are assigned by the morphology-to-performance response surface.

\subsection*{S2.2 Repeated cross-validation}

Repeated five-fold cross-validation gives mean nMAE values of 0.0176 for EUIt, 0.0265 for EG, and 0.0145 for H. The corresponding Spearman correlations are 0.991, 0.984, and 0.995. These values support low error within the analytic response-generation regime.

\subsection*{S2.3 Residual and target-tail assessment}

\begin{figure}[H]
\centering
\sifigure{A2_residual_\diagword.pdf}
\caption{Residual assessment for EUIt, EG, and H under repeated cross-validated surrogate predictions.}
\label{fig:s3-residual-assessment}
\end{figure}

The residual distributions are centered near zero for all three targets, indicating that the selected surrogate does not show a systematic offset under repeated cross-validation. The EG residuals are wider relative to their target range than the EUIt and H residuals, which is consistent with the slightly higher nMAE reported for EG in the validation summaries.

\subsection*{S2.4 Leave-one-OSLI-out assessment}

The leave-one-OSLI-out protocol with nine folds gives nMAE values of 0.0185 for EUIt, 0.0342 for EG, and 0.0185 for H. The results remain close to the repeated cross-validation regime, but EG shows the largest relative increase when an open-space-position class is withheld.

\subsection*{S2.5 Outer-shell and feature-tail holdouts}

The outer-shell holdout gives nMAE values of 0.0278 for EUIt, 0.0335 for EG, and 0.0208 for H. The feature-tail holdout gives nMAE values of 0.0230 for EUIt, 0.0287 for EG, and 0.0198 for H. These holdouts stress descriptor-domain boundaries and target tails within the analytic response-generation regime.

\section*{S3. DDPG training and benchmark comparison details}

\subsection*{S3.1 DDPG seed-level behaviour}

Each DDPG scenario uses 20 seeds, 600 episodes, and 40 surrogate interactions per episode. The seed-level summaries record best episode return, final episode return, plateau episode, and best-to-final regression. These quantities describe policy-training behaviour under serialized surrogate queries on a static surrogate response surface.

\begin{figure}[H]
\centering
\sifigure{B1_seed_\diagword.pdf}
\caption{Seed-level DDPG episode-return assessment across the three preference scenarios.}
\label{fig:s4-seed-assessment}
\end{figure}

Figure~\ref{fig:s4-seed-assessment} complements the training trajectories in the main manuscript by summarizing seed-level outcomes. The best episode return and final episode return panels distinguish the best point found during training from the final episode state. The plateau and best-to-final gap panels show that the DDPG runs can reach high-return regions while still retaining late-stage variability across seeds.

\subsection*{S3.2 Retained optimizer outputs}

DDPG retains one best scalarized candidate per seed, while NSGA-II retains final populations. CMA-ES and RandomSearch retain top-reward queried archives for secondary comparison, and FeasiblePoolRandom provides a deterministic feasible-pool reference. This retained-output structure prevents full-archive sizes from being interpreted as equivalent across optimizers.

\subsection*{S3.3 Equal-size sensitivity}

The main manuscript presents equal-size results at retained size 20 because that is the valid common retained size for DDPG and NSGA-II. Oversized equal-size requests for smaller source archives are labelled not applicable and excluded from metric interpretation.

\begin{table}[H]
\caption{Equal-size-20 descriptor-space metrics under the fixed normalized reference front.}
\label{tab:s-equal-size}
\centering
\footnotesize
\begin{tabularx}{\linewidth}{@{}p{0.28\linewidth}rrp{0.28\linewidth}@{}}
\toprule
Group & HV & IGD & Interpretation \\
\midrule
DDPG Balanced & 0.661912 & 0.465522 & weakest DDPG setting \\
DDPG Saving & 1.002757 & 0.276529 & intermediate DDPG setting \\
DDPG Generation & 1.170428 & 0.097551 & strongest DDPG setting \\
NSGA-II & 1.330995 & 0.004710 & near-ceiling retained-size result \\
CMA-ES Balanced & 1.331000 & 0.004946 & secondary corner-saturation reference \\
RandomSearch Balanced & 0.423--0.748 & 0.497--0.743 & repetition-dependent uninformed reference \\
\bottomrule
\end{tabularx}
\end{table}

The equal-size-20 comparison is the common retained-size setting available for the DDPG and NSGA-II main comparison. Larger requested sizes are not treated as comparable for methods that do not retain enough distinct candidates. For this reason, the full-archive metrics and equal-size metrics are interpreted separately in the main manuscript.

\subsection*{S3.4 HV ceiling and clipped-output behaviour}

The theoretical maximum HV under the fixed normalized reference front is 1.331. NSGA-II reaches a fraction of 1.0 without clipped-utopia rows, while CMA-ES reaches the same ceiling because nearly all retained rows collapse to clipped-utopia output tuples. This distinction is why ceiling proximity must be read together with unique tuple counts.

\begin{figure}[H]
\centering
\sifigure{B3_hv_ceiling_\diagword.pdf}
\caption{HV ceiling and clipped-output assessment for secondary optimizer archives.}
\label{fig:s5-hv-ceiling}
\end{figure}

Figure~\ref{fig:s5-hv-ceiling} explains why near-ceiling HV values must be interpreted with objective-tuple diversity. CMA-ES can reach the theoretical HV ceiling because many retained rows collapse to clipped-utopia tuples, whereas NSGA-II approaches the same reference volume with many more distinct clipped objective tuples. The figure therefore separates reference-volume coverage from archive diversity.

\subsection*{S3.5 CMA-ES and RandomSearch comparisons}

CMA-ES and RandomSearch are retained as secondary comparators. CMA-ES reaches the HV ceiling in descriptor space but does so with severe objective-tuple collapse. RandomSearch retains many unique tuples but remains weaker than NSGA-II under the main equal-size descriptor-space comparison. These results are used to interpret surrogate exploitability and uninformed-search behaviour alongside the primary DDPG and NSGA-II comparison.

\section*{S4. Morphology interpretation}

\subsection*{S4.1 Morphology descriptor signatures}

Median descriptor signatures summarize where retained outputs lie in descriptor space. NSGA-II retained outputs have median FAR = 1.71, median SVF = 0.654, and median OSLI = 6.21, while DDPG Balanced outputs have median FAR = 1.62 and median SVF = 0.702. These signatures describe retained-output tendencies and are not stable design rules.

\begin{figure}[H]
\centering
\sifigure{B2_morphology_signatures.pdf}
\caption{Median morphology descriptor signatures for representative retained-output groups.}
\label{fig:s6-morphology-signatures}
\end{figure}

The morphology signatures summarize retained-output tendencies rather than prescriptive design rules. They are useful for comparing where optimizers tend to place candidates in descriptor space, but they should be read together with the feasible-projection analysis because continuous descriptor outputs may map to repeated generated blocks.

\subsection*{S4.2 Surrogate response profiles}

\begin{figure}[H]
\centering
\sifigure{B5_nonlinear_response_profiles.pdf}
\caption{Selected surrogate response profiles for OSR, FAR, SVF, and theta.}
\label{fig:s7-response-profiles}
\end{figure}

The response profiles isolate one descriptor at a time around representative reference values. They provide a qualitative view of the local shape of the selected surrogate response surface and help interpret why certain descriptors are emphasized by the optimizer outputs. Since the profiles are surrogate-based, they are used for interpretation only and are not a replacement for physics-based calculation.

\section*{S5. Feasible projection and physics-based stress-test details}

\subsection*{S5.1 Complete projection metrics}

The complete projection table records retained rows, unique matched sample counts, projection-distance summaries, duplicate-collapse rates, and fixed-reference projected HV and IGD. NSGA-II contracts from 2000 retained descriptor candidates to 51 unique projected feasible blocks, while DDPG scenarios contract from 20 retained descriptor candidates to 17 or 18 unique generated feasible blocks.

Complete projection metrics are kept in the supplementary material because they depend on representation family, duplicate handling, and the fixed normalized reference front. The main manuscript focuses on the most interpretable projection outcomes: candidate compression and projection distance.

\subsection*{S5.2 Optimizer-linked six-case decomposition}

The physics-based stress-test set contains 24 locked cases. The 18 direct feasible cases are used to compare analytic response-generator values with physics-based outputs without optimizer projection. The six optimizer-linked cases decompose the path from descriptor-space candidate to nearest generated block and then to physics-based result. These two subsets should not be merged into a single parity plot because they answer different questions.

\begin{figure}[H]
\centering
\sifigure{B4_optimizer_linked_gap_decomposition.pdf}
\caption{Optimizer-linked six-case projection and cross-model gap decomposition.}
\label{fig:s8-optimizer-linked-gaps}
\end{figure}

Figure~\ref{fig:s8-optimizer-linked-gaps} separates the projection gap from the cross-model gap for the six optimizer-linked cases. This separation is important because an optimizer candidate is first mapped to a generated feasible block before the physics-based calculation is evaluated. The plotted gaps should therefore be interpreted as a bridge analysis between descriptor-space search and case-level physical outputs.

\subsection*{S5.3 Per-case physical results}

The 18 direct feasible cases completed both EnergyPlus and Radiance-component runs. EUIt agreement is weak, EG is represented by a GHI-based rooftop-PV proxy, and H is represented by January 20 windowsill direct-sun hours. The nMAE and rank-correlation values therefore define a stress-test boundary for the tested direct cases.

\begin{table}[H]
\caption{Direct feasible physical stress-test summary.}
\label{tab:s-physical-summary}
\centering
\footnotesize
\begin{tabularx}{\linewidth}{@{}p{0.25\linewidth}rrrrp{0.23\linewidth}@{}}
\toprule
Target & Count & MAE & nMAE & Spearman rho & Boundary \\
\midrule
EUIt & 18 & 3.788 & 0.340 & -0.404 & limited metric consistency \\
EG GHI proxy & 18 & 0.441 & 0.405 & -0.023 & rooftop-PV proxy only \\
H & 18 & 1.108 & 0.601 & -0.043 & January 20 direct-sun hours only \\
\bottomrule
\end{tabularx}
\end{table}

The physical stress-test summary shows that all direct feasible cases completed, but the agreement between analytic targets and physics-based outputs remains limited. These values define the cross-model comparison boundary used in the main manuscript. The optimizer-linked cases are not merged with this table because they include an additional projection step.

\subsection*{S5.4 EG roof-irradiance unavailability note}

The locked physical workflow did not produce annual roof-plane irradiance outputs for the direct stress-test cases. The unavailable roof-irradiance pathway is therefore kept separate from the completed EUIt, EG proxy, and H calculations.

EG is summarized with the GHI-based rooftop-PV proxy in the stress-test table. This proxy preserves a generation-related comparison across the locked cases, but it should be interpreted as a simplified rooftop-PV indicator rather than a full annual roof-irradiance result.

\section*{S6. Cross-climate case details}

\subsection*{S6.1 Weather metadata}

The baseline weather file corresponds to the station used for the case-study physical workflow. The cross-climate sensitivity analysis evaluates the same four generated feasible blocks under Harbin, Beijing, and Guangzhou weather. The climate buckets are severe cold for Harbin, cold for Beijing, and hot-summer warm-winter for Guangzhou. The TMY period is 2009--2023 for these weather files.

\subsection*{S6.2 Four-block climate results}

The climate sensitivity set uses the same four block morphologies under three additional EPW files. The surrogate remains fixed, and the morphology set is held constant after the climate results are computed. The results therefore describe weather sensitivity for selected block morphologies within the tested climate set.

\begin{table}[H]
\caption{Station-level mean physical changes relative to baseline weather.}
\label{tab:s-climate-summary}
\centering
\small
\setlength{\tabcolsep}{5pt}
\renewcommand{\arraystretch}{1.08}
\begin{threeparttable}
\begin{tabular*}{\linewidth}{@{\extracolsep{\fill}}lrrr>{\raggedright\arraybackslash}p{0.34\linewidth}@{}}
\toprule
Station & $\Delta$EUIt & $\Delta$EG & $\Delta$H & Main interpretation \\
\midrule
Beijing & 13.467 & 0.135 & -0.056 & higher demand under the cold-climate weather file \\
Guangzhou & -2.972 & -0.057 & -0.025 & lower demand under the hot-summer warm-winter weather file \\
Harbin & 55.681 & -0.044 & -0.675 & strong demand increase under the severe-cold weather file \\
\bottomrule
\end{tabular*}
\begin{tablenotes}[flushleft]
\footnotesize
\item $\Delta$EUIt is in kWh m$^{-2}$ y$^{-1}$, $\Delta$EG is in $10^6$ kWh y$^{-1}$, and $\Delta$H is in h. Positive values denote increases relative to the baseline weather file.
\end{tablenotes}
\end{threeparttable}
\end{table}

Table~\ref{tab:s-climate-summary} reports mean target changes by weather file, while Fig.~\ref{fig:s9-climate-detail} shows the corresponding case-level variation across the four selected blocks. The results indicate that EUIt is strongly climate-sensitive, especially under the severe-cold weather file, whereas EG changes are smaller and preserve the ordering of the selected blocks.

\begin{figure}[H]
\centering
\sifigure{B6_climate_case_detail.pdf}
\caption{Case-level climate sensitivity for EUIt, EG, and H across Beijing, Guangzhou, and Harbin.}
\label{fig:s9-climate-detail}
\end{figure}

\subsection*{S6.3 Rank stability across climates}

The rank-stability results indicate target-specific climate sensitivity. EG has stable ordering across the three additional climates, with Spearman $\rho$ equal to 1.0 in each case. H shows lower ordering stability, with Spearman $\rho$ values of 0.2 for Beijing, 0.4 for Guangzhou, and 0.2 for Harbin. EUIt ordering remains stable for Beijing and Guangzhou but changes under Harbin, where Spearman $\rho$ is 0.2. These results support a limited climate-sensitivity interpretation for the selected four-block set.

\end{document}
